\IfFileExists{stacks-project.cls}{%
\documentclass{stacks-project}
}{%
\documentclass{amsart}
}

% For dealing with references we use the comment environment
\usepackage{verbatim}
\newenvironment{reference}{\comment}{\endcomment}
%\newenvironment{reference}{}{}
\newenvironment{slogan}{\comment}{\endcomment}
\newenvironment{history}{\comment}{\endcomment}

% For commutative diagrams we use Xy-pic
\usepackage[all]{xy}

% We use 2cell for 2-commutative diagrams.
\xyoption{2cell}
\UseAllTwocells

% We use multicol for the list of chapters between chapters
\usepackage{multicol}

% This is generally recommended for better output
\usepackage{lmodern}
\usepackage[T1]{fontenc}

% For cross-file-references
\usepackage{xr-hyper}

% Package for hypertext links:
\usepackage{hyperref}

% For any local file, say "hello.tex" you want to link to please
% use \externaldocument[hello-]{hello}
\externaldocument[introduction-]{introduction}
\externaldocument[conventions-]{conventions}
\externaldocument[sets-]{sets}
\externaldocument[categories-]{categories}
\externaldocument[topology-]{topology}
\externaldocument[sheaves-]{sheaves}
\externaldocument[sites-]{sites}
\externaldocument[stacks-]{stacks}
\externaldocument[fields-]{fields}
\externaldocument[algebra-]{algebra}
\externaldocument[brauer-]{brauer}
\externaldocument[homology-]{homology}
\externaldocument[derived-]{derived}
\externaldocument[simplicial-]{simplicial}
\externaldocument[more-algebra-]{more-algebra}
\externaldocument[smoothing-]{smoothing}
\externaldocument[modules-]{modules}
\externaldocument[sites-modules-]{sites-modules}
\externaldocument[injectives-]{injectives}
\externaldocument[cohomology-]{cohomology}
\externaldocument[sites-cohomology-]{sites-cohomology}
\externaldocument[dga-]{dga}
\externaldocument[dpa-]{dpa}
\externaldocument[sdga-]{sdga}
\externaldocument[hypercovering-]{hypercovering}
\externaldocument[schemes-]{schemes}
\externaldocument[constructions-]{constructions}
\externaldocument[properties-]{properties}
\externaldocument[morphisms-]{morphisms}
\externaldocument[coherent-]{coherent}
\externaldocument[divisors-]{divisors}
\externaldocument[limits-]{limits}
\externaldocument[varieties-]{varieties}
\externaldocument[topologies-]{topologies}
\externaldocument[descent-]{descent}
\externaldocument[perfect-]{perfect}
\externaldocument[more-morphisms-]{more-morphisms}
\externaldocument[flat-]{flat}
\externaldocument[groupoids-]{groupoids}
\externaldocument[more-groupoids-]{more-groupoids}
\externaldocument[etale-]{etale}
\externaldocument[chow-]{chow}
\externaldocument[intersection-]{intersection}
\externaldocument[pic-]{pic}
\externaldocument[weil-]{weil}
\externaldocument[adequate-]{adequate}
\externaldocument[dualizing-]{dualizing}
\externaldocument[duality-]{duality}
\externaldocument[discriminant-]{discriminant}
\externaldocument[derham-]{derham}
\externaldocument[local-cohomology-]{local-cohomology}
\externaldocument[algebraization-]{algebraization}
\externaldocument[curves-]{curves}
\externaldocument[resolve-]{resolve}
\externaldocument[models-]{models}
\externaldocument[functors-]{functors}
\externaldocument[equiv-]{equiv}
\externaldocument[pione-]{pione}
\externaldocument[etale-cohomology-]{etale-cohomology}
\externaldocument[proetale-]{proetale}
\externaldocument[relative-cycles-]{relative-cycles}
\externaldocument[more-etale-]{more-etale}
\externaldocument[trace-]{trace}
\externaldocument[crystalline-]{crystalline}
\externaldocument[spaces-]{spaces}
\externaldocument[spaces-properties-]{spaces-properties}
\externaldocument[spaces-morphisms-]{spaces-morphisms}
\externaldocument[decent-spaces-]{decent-spaces}
\externaldocument[spaces-cohomology-]{spaces-cohomology}
\externaldocument[spaces-limits-]{spaces-limits}
\externaldocument[spaces-divisors-]{spaces-divisors}
\externaldocument[spaces-over-fields-]{spaces-over-fields}
\externaldocument[spaces-topologies-]{spaces-topologies}
\externaldocument[spaces-descent-]{spaces-descent}
\externaldocument[spaces-perfect-]{spaces-perfect}
\externaldocument[spaces-more-morphisms-]{spaces-more-morphisms}
\externaldocument[spaces-flat-]{spaces-flat}
\externaldocument[spaces-groupoids-]{spaces-groupoids}
\externaldocument[spaces-more-groupoids-]{spaces-more-groupoids}
\externaldocument[bootstrap-]{bootstrap}
\externaldocument[spaces-pushouts-]{spaces-pushouts}
\externaldocument[spaces-chow-]{spaces-chow}
\externaldocument[groupoids-quotients-]{groupoids-quotients}
\externaldocument[spaces-more-cohomology-]{spaces-more-cohomology}
\externaldocument[spaces-simplicial-]{spaces-simplicial}
\externaldocument[spaces-duality-]{spaces-duality}
\externaldocument[formal-spaces-]{formal-spaces}
\externaldocument[restricted-]{restricted}
\externaldocument[spaces-resolve-]{spaces-resolve}
\externaldocument[formal-defos-]{formal-defos}
\externaldocument[defos-]{defos}
\externaldocument[cotangent-]{cotangent}
\externaldocument[examples-defos-]{examples-defos}
\externaldocument[algebraic-]{algebraic}
\externaldocument[examples-stacks-]{examples-stacks}
\externaldocument[stacks-sheaves-]{stacks-sheaves}
\externaldocument[criteria-]{criteria}
\externaldocument[artin-]{artin}
\externaldocument[quot-]{quot}
\externaldocument[stacks-properties-]{stacks-properties}
\externaldocument[stacks-morphisms-]{stacks-morphisms}
\externaldocument[stacks-limits-]{stacks-limits}
\externaldocument[stacks-cohomology-]{stacks-cohomology}
\externaldocument[stacks-perfect-]{stacks-perfect}
\externaldocument[stacks-introduction-]{stacks-introduction}
\externaldocument[stacks-more-morphisms-]{stacks-more-morphisms}
\externaldocument[stacks-geometry-]{stacks-geometry}
\externaldocument[moduli-]{moduli}
\externaldocument[moduli-curves-]{moduli-curves}
\externaldocument[examples-]{examples}
\externaldocument[exercises-]{exercises}
\externaldocument[guide-]{guide}
\externaldocument[desirables-]{desirables}
\externaldocument[coding-]{coding}
\externaldocument[obsolete-]{obsolete}
\externaldocument[fdl-]{fdl}
\externaldocument[index-]{index}

% Theorem environments.
%
\theoremstyle{plain}
\newtheorem{theorem}[subsection]{Theorem}
\newtheorem{proposition}[subsection]{Proposition}
\newtheorem{lemma}[subsection]{Lemma}

\theoremstyle{definition}
\newtheorem{definition}[subsection]{Definition}
\newtheorem{example}[subsection]{Example}
\newtheorem{exercise}[subsection]{Exercise}
\newtheorem{situation}[subsection]{Situation}

\theoremstyle{remark}
\newtheorem{remark}[subsection]{Remark}
\newtheorem{remarks}[subsection]{Remarks}

\numberwithin{equation}{subsection}

% Macros
%
\def\lim{\mathop{\mathrm{lim}}\nolimits}
\def\colim{\mathop{\mathrm{colim}}\nolimits}
\def\Spec{\mathop{\mathrm{Spec}}}
\def\Hom{\mathop{\mathrm{Hom}}\nolimits}
\def\Ext{\mathop{\mathrm{Ext}}\nolimits}
\def\SheafHom{\mathop{\mathcal{H}\!\mathit{om}}\nolimits}
\def\SheafExt{\mathop{\mathcal{E}\!\mathit{xt}}\nolimits}
\def\Sch{\mathit{Sch}}
\def\Mor{\mathop{\mathrm{Mor}}\nolimits}
\def\Ob{\mathop{\mathrm{Ob}}\nolimits}
\def\Sh{\mathop{\mathit{Sh}}\nolimits}
\def\NL{\mathop{N\!L}\nolimits}
\def\CH{\mathop{\mathrm{CH}}\nolimits}
\def\proetale{{pro\text{-}\acute{e}tale}}
\def\etale{{\acute{e}tale}}
\def\QCoh{\mathit{QCoh}}
\def\Ker{\mathop{\mathrm{Ker}}}
\def\Im{\mathop{\mathrm{Im}}}
\def\Coker{\mathop{\mathrm{Coker}}}
\def\Coim{\mathop{\mathrm{Coim}}}

% Boxtimes
%
\DeclareMathSymbol{\boxtimes}{\mathbin}{AMSa}{"02}

%
% Macros for moduli stacks/spaces
%
\def\QCohstack{\mathcal{QC}\!\mathit{oh}}
\def\Cohstack{\mathcal{C}\!\mathit{oh}}
\def\Spacesstack{\mathcal{S}\!\mathit{paces}}
\def\Quotfunctor{\mathrm{Quot}}
\def\Hilbfunctor{\mathrm{Hilb}}
\def\Curvesstack{\mathcal{C}\!\mathit{urves}}
\def\Polarizedstack{\mathcal{P}\!\mathit{olarized}}
\def\Complexesstack{\mathcal{C}\!\mathit{omplexes}}
% \Pic is the operator that assigns to X its picard group, usage \Pic(X)
% \Picardstack_{X/B} denotes the Picard stack of X over B
% \Picardfunctor_{X/B} denotes the Picard functor of X over B
\def\Pic{\mathop{\mathrm{Pic}}\nolimits}
\def\Picardstack{\mathcal{P}\!\mathit{ic}}
\def\Picardfunctor{\mathrm{Pic}}
\def\Deformationcategory{\mathcal{D}\!\mathit{ef}}
\usepackage{cleveref}
\crefname{originalreference}{original reference}{original references}
\Crefname{originalreference}{Original reference}{Original references}
\makeatletter
\newlabel{spaces-lemma-base-change-representable-transformations}{{02WB}{source}{}{sourcepointer-1}{}}
\newlabel{spaces-lemma-base-change-representable-transformations@cref}{{[originalreference][0][]02WB}{[1][0][]source}{}{}{}}
\makeatother
\crefname{originalreference}{original reference}{original references}
\Crefname{originalreference}{Original reference}{Original references}
\makeatletter
\newlabel{spaces-lemma-base-change-representable-transformations-property}{{02WL}{source}{}{sourcepointer-2}{}}
\newlabel{spaces-lemma-base-change-representable-transformations-property@cref}{{[originalreference][0][]02WL}{[1][0][]source}{}{}{}}
\makeatother
\crefname{originalreference}{original reference}{original references}
\Crefname{originalreference}{Original reference}{Original references}
\makeatletter
\newlabel{spaces-lemma-etale-locally-representable-by-space-gives-space}{{0BGR}{source}{}{sourcepointer-3}{}}
\newlabel{spaces-lemma-etale-locally-representable-by-space-gives-space@cref}{{[originalreference][0][]0BGR}{[1][0][]source}{}{}{}}
\makeatother
\crefname{originalreference}{original reference}{original references}
\Crefname{originalreference}{Original reference}{Original references}
\makeatletter
\newlabel{spaces-morphisms-lemma-etale-locally-quasi-finite}{{03XU}{source}{}{sourcepointer-4}{}}
\newlabel{spaces-morphisms-lemma-etale-locally-quasi-finite@cref}{{[originalreference][0][]03XU}{[1][0][]source}{}{}{}}
\makeatother
\crefname{originalreference}{original reference}{original references}
\Crefname{originalreference}{Original reference}{Original references}
\makeatletter
\newlabel{spaces-morphisms-lemma-locally-quasi-finite-separated-representable}{{0418}{source}{}{sourcepointer-5}{}}
\newlabel{spaces-morphisms-lemma-locally-quasi-finite-separated-representable@cref}{{[originalreference][0][]0418}{[1][0][]source}{}{}{}}
\makeatother
\title{Excision equivalence for modifications over a separated etale neighbourhood}
\author{The Stacks Project Authors}
\date{stacks-project:spaces-limits.tex:local-snapshot-20261012; GNU FDL 1.2 or later}
\begin{document}
\maketitle
\noindent\textit{Editorial note.} Counted span is the author's own proof environment, source lines 4362-4417, following the statement at 4336-4360.\par\medskip
\noindent\textit{Editorial note.} No omissions; the uniqueness of the restriction of y over U is justified in the text.\par\medskip
\noindent\textit{Allowed standard prerequisites (editorial).}\begin{itemize}
\item Separated etale morphisms of spaces are representable. Morphisms of Spaces, Lemmas spaces-morphisms-lemma-locally-quasi-finite-separated-representable and spaces-morphisms-lemma-etale-locally-quasi-finite (cross-chapter, cited at \textasciitilde{}4400-4404). f : V -> W is representable.
\item Base change of representable transformations. Spaces, Lemmas spaces-lemma-base-change-representable-transformations and spaces-lemma-base-change-representable-transformations-property (cross-chapter, cited at \textasciitilde{}4411-4413). W\textasciicircum{}0 amalga Y -> X is representable, surjective, etale.
\item Criterion for being an algebraic space. Spaces, Lemma spaces-lemma-etale-locally-representable-by-space-gives-space (cross-chapter, cited at \textasciitilde{}4414-4416). X is an algebraic space.
\end{itemize}\medskip
\ifdefined\hypertarget\hypertarget{sourcepointer-1}{}\fi
\noindent\textit{Source reference.} Original reference 02WB: focused-seeds/geometry/S9-spaces-limits-excision-modifications/ref-tags.txt, source lines 5--5. Navigation pointer only: the counted proof cites this statement by \textbackslash{}ref as a standing input; the deduction the counted proof has to make is carried out inside the counted span, and no part of that deduction is delegated to this pointer.\par\medskip
\ifdefined\hypertarget\hypertarget{sourcepointer-2}{}\fi
\noindent\textit{Source reference.} Original reference 02WL: focused-seeds/geometry/S9-spaces-limits-excision-modifications/ref-tags.txt, source lines 6--6. Navigation pointer only: the counted proof cites this statement by \textbackslash{}ref as a standing input; the deduction the counted proof has to make is carried out inside the counted span, and no part of that deduction is delegated to this pointer.\par\medskip
\ifdefined\hypertarget\hypertarget{sourcepointer-3}{}\fi
\noindent\textit{Source reference.} Original reference 0BGR: focused-seeds/geometry/S9-spaces-limits-excision-modifications/ref-tags.txt, source lines 7--7. Navigation pointer only: the counted proof cites this statement by \textbackslash{}ref as a standing input; the deduction the counted proof has to make is carried out inside the counted span, and no part of that deduction is delegated to this pointer.\par\medskip
\ifdefined\hypertarget\hypertarget{sourcepointer-4}{}\fi
\noindent\textit{Source reference.} Original reference 03XU: focused-seeds/geometry/S9-spaces-limits-excision-modifications/ref-tags.txt, source lines 8--8. Navigation pointer only: the counted proof cites this statement by \textbackslash{}ref as a standing input; the deduction the counted proof has to make is carried out inside the counted span, and no part of that deduction is delegated to this pointer.\par\medskip
\ifdefined\hypertarget\hypertarget{sourcepointer-5}{}\fi
\noindent\textit{Source reference.} Original reference 0418: focused-seeds/geometry/S9-spaces-limits-excision-modifications/ref-tags.txt, source lines 9--9. Navigation pointer only: the counted proof cites this statement by \textbackslash{}ref as a standing input; the deduction the counted proof has to make is carried out inside the counted span, and no part of that deduction is delegated to this pointer.\par\medskip
\begin{lemma}
\label{lemma-excision-modifications}
Let $S$ be a scheme. Consider a separated \'etale morphism
$f : V \to W$ of algebraic spaces over $S$.
Assume there exists a
closed subspace $T \subset W$ such that $f^{-1}T \to T$ is
an isomorphism. Then, with $W^0 = W \setminus T$ and
$V^0 = f^{-1}W^0$ the base change functor
$$
\left\{
\begin{matrix}
g : X \to W\text{ morphism of algebraic spaces} \\
g^{-1}(W^0) \to W^0\text{ is an isomorphism}
\end{matrix}
\right\}
\longrightarrow
\left\{
\begin{matrix}
h : Y \to V\text{ morphism of algebraic spaces} \\
h^{-1}(V^0) \to V^0\text{ is an isomorphism}
\end{matrix}
\right\}
$$
is an equivalence of categories.
\end{lemma}

\begin{proof}
Since $V \to W$ is separated we see that
$V \times_W V = \Delta(V) \amalg U$ for some open and closed subspace
$U$ of $V \times_W V$. By the assumption that $f^{-1}T \to T$ is an
isomorphism we see that $U \times_W T = \emptyset$, i.e., the two
projections $U \to V$ maps into $V^0$.

\medskip\noindent
Given $h : Y \to V$ in the right hand category, consider the
contravariant functor $X$ on $(\Sch/S)_{fppf}$ defined by the rule
$$
X(T) = \{(w, y) \mid
w :  T \to W,\ y : T \times_{w, W} V \to Y\text{ morphism over }V\}
$$
Denote $g : X \to W$ the map sending $(w, y) \in X(T)$ to $w \in W(T)$.
Since $h^{-1}V^0 \to V^0$ is an isomorphism, we see that if
$w : T \to W$ maps into $W^0$, then there is a unique choice for $h$.
In other words $X \times_{g, W} W^0 = W^0$. On the other hand, consider
a $T$-valued point $(w, y, v)$ of $X \times_{g, W, f} V$.
Then $w = f \circ v$ and
$$
y : T \times_{f \circ v, W} V \longrightarrow V
$$
is a morphism over $V$. Consider the morphism
$$
T \times_{f \circ v, W} V \xrightarrow{(v, \text{id}_V)}
V \times_W V = V \amalg U
$$
The inverse image of $V$ is $T$ embedded via
$(\text{id}_T, v) : T \to T \times_{f \circ v, W} V$.
The composition $y' = y \circ (\text{id}_T, v) : T \to Y$
is a morphism with $v = h \circ y'$ which determines $y$ because the
restriction of $y$ to the other part is uniquely determined as
$U$ maps into $V^0$ by the second projection. It follows that
$X \times_{g, W, f} V \to Y$, $(w, y, v) \mapsto y'$ is an isomorphism.

\medskip\noindent
Thus if we can show that $X$ is an algebraic space, then we are done.
Since $V \to W$ is separated and \'etale it is representable by
Morphisms of Spaces, Lemma
\ref{spaces-morphisms-lemma-locally-quasi-finite-separated-representable}
(and Morphisms of Spaces, Lemma
\ref{spaces-morphisms-lemma-etale-locally-quasi-finite}).
Of course $W^0 \to W$ is representable and \'etale as it is an
open immersion. Thus
$$
W^0 \amalg Y = X \times_{g, W} W^0 \amalg X \times_{g, W, f} V
= X \times_{g, W} (W^0 \amalg V) \longrightarrow X
$$
is representable, surjective, and \'etale by Spaces, Lemmas
\ref{spaces-lemma-base-change-representable-transformations} and
\ref{spaces-lemma-base-change-representable-transformations-property}.
Thus $X$ is an algebraic
space by Spaces, Lemma
\ref{spaces-lemma-etale-locally-representable-by-space-gives-space}.
\end{proof}

\end{document}
